\par\Needspace{6\baselineskip}\medskip
\begin{center}\hypertarget{native-199-article149}{}\pdfbookmark[2]{149.}{article-149}149.\end{center}
In ultima demum hypothesi elementa ipsa calculabuntur, vel ex $f,r',r''$,
vel ex $f'',r,r'$, perducendo scilicet ad finem calculum alterutrum, quem in hypothesibus
antecedentibus tantummodo usque ad $\eta$ vel $\eta''$ prosequi oportuerat:
si utrumque perficere placuerit, harmonia numerorum resultantium novam totius
laboris confirmationem suppeditabit. Attamen praestat, quam primum $f,f',f''$
erutae sunt, elementa e sola combinatione loci primi cum tertio derivare,
puta ex $f',r,r''$ atque temporis intervallo, tandemque ad maiorem calculi
certitudinem locum medium in orbita secundum elementa inventa determinare.

Hoc itaque modo sectionis conicae dimensiones innotescent, puta excentricitas,
semiaxis maior sive semiparameter, positio perihelii respectu locorum
heliocentricorum $C,C',C''$, motus medius, atque anomalia media pro epocha
arbitraria, siquidem orbita elliptica est, vel tempus transitus per perihelium,
si orbita fit hyperbolica vel parabolica. Superest itaque tantummodo, ut positio
locorum heliocentricorum in orbita respectu nodi ascendentis, positio huius
nodi respectu puncti aequinoctialis, atque inclinatio orbitae ad eclipticam (vel
aequatorem) determinentur. Haec omnia per solutionem unius trianguli sphaerici
efficere licet. Sit $\AscendingNode$ longitudo nodi ascendentis; $i$ inclinatio orbitae; $g$ et
$g''$ argumenta latitudinis in observatione prima et tertia; denique $l-\AscendingNode=h$,
$l''-\AscendingNode=h''$. Exprimente iam in fig. quarta $\AscendingNode$ nodum ascendentem, trianguli
$\AscendingNode AC$ latera erunt $AD'-\zeta,g,h$, angulique his resp. oppositi $i,180^{0}-\gamma,u$.
Habebimus itaque
\SourceLeaf{200}
\[
\begin{aligned}
\sin\tfrac12i\sin\tfrac12(g+h)&=\sin\tfrac12(AD'-\zeta)\sin\tfrac12(\gamma+u)\\
\sin\tfrac12i\cos\tfrac12(g+h)&=\cos\tfrac12(AD'-\zeta)\sin\tfrac12(\gamma-u)\\
\cos\tfrac12i\sin\tfrac12(g-h)&=\sin\tfrac12(AD'-\zeta)\cos\tfrac12(\gamma+u)\\
\cos\tfrac12i\cos\tfrac12(g-h)&=\cos\tfrac12(AD'-\zeta)\cos\tfrac12(\gamma-u).
\end{aligned}
\]

Duae primae aequationes dabunt $\tfrac12(g+h)$ et $\sin\tfrac12i$, duae reliquae $\tfrac12(g-h)$
et $\cos\tfrac12i$; ex $g$ innotescet situs perihelii respectu nodi ascendentis, ex $h$ situs
nodi in ecliptica; denique innotescet $i$, sinu et cosinu se mutuo confirmantibus.
Ad eundem scopum pervenire possumus adiumento trianguli $\AscendingNode A''C''$,
ubi tantummodo in formulis praecedentibus characteres $g,h,A,\zeta,\gamma,u$ in $g'',h'',A'',\zeta'',\gamma'',u''$
mutare oportet. Ut toti labori adhuc alia confirmatio concilietur,
haud abs re erit, calculum utroque modo perficere: unde si quae leviusculae
differentiae inter valores ipsius $i,\AscendingNode$ atque longitudinis perihelii in
orbita prodeunt, valores medios adoptare conveniet. Raro tamen hae differentiae
ad $0''{,}1$ vel $0''{,}2$ ascendent, siquidem omnes calculi septem figuris decimalibus
accurate elaborati fuerant.

Ceterum quoties loco eclipticae aequator tamquam planum fundamentale
adoptatum est, nulla hinc in calculo differentia orietur, nisi quod loco punctorum
$A,A''$ intersectiones aequatoris cum circulis maximis $AB,A''B''$ accipiendae
sunt.
